\documentclass[11pt]{article}
\usepackage[a4paper,margin=25mm]{geometry}
\usepackage{amsmath,amssymb,amsthm,booktabs,array}
\usepackage[T1]{fontenc}
\usepackage{lmodern}
\usepackage[hidelinks]{hyperref}
\newtheorem{theorem}{Theorem}
\newcommand{\N}{\mathbb N}
\newcommand{\R}{\mathbb R}
\title{A proof of conjecture 00000000106}
\author{C0ldSmi1e}
\date{October 5, 2026}
\begin{document}
\maketitle
\begin{abstract}
For each integer $k\ge2$, the exponent $c_k=1/(2k)$ works: there are
infinitely many natural numbers $n$ for which every one of
$n+1,\ldots,n+k$ has a prime factor larger than $n^{c_k}$.
We prove this using Bertrand's postulate and a bounded Chinese remainder
representative. The accompanying Lean project proves the assertion for the
actual largest prime factor, with a single exponent fixed for each $k$.
\end{abstract}

\section{Statement and interpretation}
The English source states:
\begin{quote}
Conjecture: (Family of consecutive rough numbers) For every $k\ge2$ there
exists $c_k>0$ such that infinitely many $n$ satisfy: the largest prime factor
of each of $n+1,\ldots,n+k$ exceeds $n^{c_k}$.
\end{quote}
The Chinese source gives the same quantifiers and largest-prime-factor
condition. Let $P(m)$ denote the greatest prime divisor of $m\ge2$.
The explicit condition concerns the \emph{largest} prime factor; the word
``rough'' in the title introduces no additional condition on smaller factors.
All powers in the statement are real powers.

We prove a stronger, explicit formulation with $n\ge2$:
\begin{theorem}\label{thm:main}
For every $k\in\N$ with $k\ge2$ and every $N\in\N$, there is an integer
$n>N$, with $n\ge2$, such that
\[
  n^{1/(2k)}<P(n+i)\qquad(1\le i\le k).
\]
Consequently the conjecture holds with $c_k=1/(2k)>0$.
\end{theorem}
Restricting to $n\ge2$ suffices for the original infinitude claim, and ensures
that every argument of $P$ is at least two. The exponent depends only on $k$;
neither the bound $N$ nor the offset $i$ changes it.

\section{Proof}
Fix $k\ge2$ and $N\in\N$, and put
\[
 t=\max\{2^k,N+k+3\}.
\]
Then $t>k$, $t>N$, $t\ge2$, and $2^k\le t$.
By Bertrand's postulate, for each $j=0,\ldots,k-1$ choose a prime $p_j$ with
\begin{equation}\label{eq:primes}
  2^j t<p_j\le 2^{j+1}t.
\end{equation}
These primes are strictly increasing: if $j<\ell$, then
$p_j\le2^{j+1}t\le2^\ell t<p_\ell$. Hence they are pairwise coprime.
Also $p_j>t>k\ge j+1$, so $p_j-(j+1)$ is a natural number.

The Chinese remainder theorem supplies the representative $n$ satisfying
\begin{equation}\label{eq:crt}
  0\le n<\prod_{j=0}^{k-1}p_j,
  \qquad n\equiv p_j-(j+1)\pmod{p_j}
  \quad(0\le j<k).
\end{equation}
It follows that $p_j\mid n+(j+1)$ for every $j$.
In particular, $p_0\mid n+1$. Because $n+1$ is positive,
$p_0\le n+1$. The strict integer inequality $t<p_0$ then implies
$t\le n$. This lower bound is essential: the bounded CRT representative
grows beyond every prescribed $N$, since $n\ge t>N$.

For every $j<k$, \eqref{eq:primes} and $2^k\le t$ give
\[
  p_j\le2^{j+1}t\le2^kt\le t^2.
\]
Combining this with the strict bound in \eqref{eq:crt} yields
\[
  n<\prod_{j=0}^{k-1}p_j\le(t^2)^k=t^{2k}.
\]
Since $n,t$ are positive and $2k>0$, monotonicity of positive real powers gives
\[
  n^{1/(2k)}<t<p_j\le P(n+j+1)
  \qquad(0\le j<k).
\]
Substituting $i=j+1$ proves the simultaneous inequalities in
Theorem~\ref{thm:main}. The construction works for every $N$, so the set of
qualifying natural numbers is unbounded and therefore infinite.
\hfill$\square$

\section{Formalization and verification}
The project uses Lean 4.19.0 and Mathlib v4.19.0. Its main file is
\texttt{Conjecture106.lean}, in namespace \texttt{Conjecture106}.
For $m\ge2$, its definition reads
\[
  \texttt{largestPrimeFactor}(m)
    =\sup\{p\in\N:p\text{ is a prime divisor of }m\}
\]
using the finite set \texttt{Nat.primeFactors}. At zero and one this
finite-set convention returns zero. Those inputs are never used in the
conjecture. The theorem \texttt{largestPrimeFactor\_spec} proves, for every
$m\ge2$, that the defined value is prime, divides $m$, and is at least every
prime divisor of $m$. The theorem \texttt{threshold\_iff\_prime\_divisor}
proves the exact equivalence
\[
  x<P(m)\quad\Longleftrightarrow\quad
  \exists p\in\N\ (p\text{ prime},\ p\mid m,\ x<p)
  \qquad(m\ge2, x\in\R).
\]
Thus the chosen prime divisors establish precisely the asserted condition
on the greatest prime divisor.

\noindent The remaining declarations follow the mathematical proof:
\begin{center}
\begin{tabular}{@{}>{\raggedright\arraybackslash}p{0.37\linewidth}
                    >{\raggedright\arraybackslash}p{0.57\linewidth}@{}}
\toprule
Declaration & Meaning \\
\midrule
\texttt{Good}, \texttt{Statement} & The simultaneous strict inequalities
and the full quantifier order of the conjecture. \\
\texttt{dyadic\_primes} & Strictly increasing primes in successive intervals
$(2^j t,2^{j+1}t]$. \\
\texttt{bounded\_representative} & A CRT representative with
$t\le n<t^{2k}$ and the required prime divisors. \\
\texttt{unbounded\_good} & A qualifying $n>N$ for every bound $N$,
with the fixed exponent $1/(2k)$. \\
\texttt{conjecture} & The full statement, including ordinary set infinitude. \\
\bottomrule
\end{tabular}
\end{center}

The project pins all dependencies in \texttt{lake-manifest.json}.
The included verification records distinguish compilation and local semantic
review from any subsequent maintainer decision. They record a fresh build
of the submitted sources against the pinned dependency cache, replay with
warnings treated as errors, declaration types and axiom dependencies, and
inspection of all constants originating in the submitted proof module.
There are no admitted proofs, custom axioms, or \texttt{native\_decide} calls.
The accepted axiom dependencies are only Lean's standard
\texttt{propext}, \texttt{Classical.choice}, and \texttt{Quot.sound}.
The mathematical argument requires no external numerical computation.
Auxiliary inspection programs and their execution records are supplied for
reproducibility; they are not imported by the proof.

\paragraph{Main verification commands.}
From the included Lean project, after installing its pinned dependencies:
\begin{verbatim}
lake build
lake env lean -DwarningAsError=true Conjecture106.lean
lake env lean -DwarningAsError=true Check.lean
\end{verbatim}
The submission's verification instructions explain how to rerun the full
auxiliary audit and locate its recorded results.

\paragraph{Library results.}
The construction uses Mathlib's
\texttt{Nat.exists\_prime\_lt\_and\_le\_two\_mul},
\texttt{Nat.chineseRemainderOfFinset}, and
\texttt{Nat.chineseRemainderOfFinset\_lt\_prod}.
The conversion to the real exponent uses
\texttt{Real.rpow\_inv\_lt\_iff\_of\_pos} and
\texttt{Real.rpow\_natCast}; infinitude uses
\texttt{Set.infinite\_iff\_exists\_gt} on the natural numbers.
\end{document}
